% !TeX program = lualatex
% Standalone research notebook for original catalogue problem 66.
% Compile twice: lualatex 66.tex
% Original entry retained; research subsection and [E] references added.
% No external bibliography, image, data, or style file is required.
\documentclass[11pt,a4paper]{article}
\usepackage[margin=24mm,headheight=14pt]{geometry}
\usepackage{fontspec}
\setmainfont{Latin Modern Roman}
\setsansfont{Latin Modern Sans}
\setmonofont{Latin Modern Mono}
\usepackage{amsmath,amssymb,mathtools,amsthm}
\usepackage{unicode-math}
\setmathfont{Latin Modern Math}
\usepackage{microtype}
\usepackage{enumitem,xurl,needspace}
\usepackage[dvipsnames]{xcolor}
\usepackage{hyperref}
\hypersetup{unicode=true,colorlinks=true,linkcolor=MidnightBlue,
 urlcolor=MidnightBlue,bookmarksnumbered=true,
 pdftitle={Research notebook - Problem 66},
 pdfauthor={Research attempt; not independently refereed}}
\usepackage{bookmark}
\usepackage{titlesec}
\titleformat{\section}{\Large\bfseries\raggedright}{\thesection}{0.7em}{}
\usepackage{fancyhdr}
\pagestyle{fancy}\fancyhf{}
\fancyhead[L]{\small\sffamily Research notebook}
\fancyhead[R]{\small\sffamily Problem 66}
\fancyfoot[L]{\scriptsize 27 September 2026}
\fancyfoot[C]{\thepage}
\fancyfoot[R]{\scriptsize Unrefereed research attempt}
\setlength{\parindent}{0pt}
\setlength{\parskip}{4pt plus 1pt}
\setlength{\emergencystretch}{3em}
\clubpenalty=10000
\widowpenalty=10000
\displaywidowpenalty=10000
\setcounter{secnumdepth}{1}
\setlist[itemize]{leftmargin=3em,itemsep=2pt,topsep=3pt,parsep=0pt,before=\small}
\allowdisplaybreaks
\newcommand{\N}{\mathbb{N}}
\newcommand{\Z}{\mathbb{Z}}
\newcommand{\Q}{\mathbb{Q}}
\newcommand{\R}{\mathbb{R}}
\newcommand{\C}{\mathbb{C}}
\newcommand{\F}{\mathbb{F}}
\newcommand{\A}{\mathbb{A}}
\newcommand{\PP}{\mathbb P}
\newcommand{\E}{\mathbb{E}}
\newcommand{\eps}{\varepsilon}
\newcommand{\dd}{\,\mathrm d}
\newcommand{\sref}[2]{\hyperlink{source:#1:#2}{[S#2]}}
\newcommand{\eref}[2]{\hyperlink{extra:#1:#2}{[E#2]}}
\newif\ifshowcatalogue
\showcataloguetrue
\newcommand{\cataloguescope}[1]{\ifshowcatalogue
 \paragraph{Exact scope in the supplied catalogue.}
 \begin{quote}\small #1\end{quote}\fi}
\newtheorem{notebooklemma}{Lemma}[section]
\newtheorem{notebookproposition}[notebooklemma]{Proposition}
\numberwithin{equation}{section}
\begin{document}
\setcounter{section}{65}
% ================================================================
% problemId: problem.artins-conjecture-on-the-holomorphy-of-non-abelian-artin-l-functions
% source releaseRank: 66; source releaseStatus: open
% exactTarget SHA-256: 2ca93fe3bf7718ff09d28eab43035585a52f208a0dbefe1e424a8e11b62bf81f
\clearpage
\section{\texorpdfstring{Artins Conjecture on the Holomorphy of Non-Abelian Artin L-Functions}{Artins Conjecture on the Holomorphy of Non-Abelian Artin L-Functions}}\label{66}
\subsection{Definitions and mathematical statement}
An Artin representation is a continuous complex representation $\rho:G_K\to\operatorname{GL}(V)$ with finite image. Its Euler product has local factors
\[
 L_v(s,\rho)=\det(1-\rho(\operatorname{Frob}_v)(Nv)^{-s}\mid V^{I_v})^{-1},
\]
with compatible Frobenius convention and inertia invariants. For every irreducible nontrivial such $\rho$, the conjecture says that the resulting meromorphic continuation $L(s,\rho)$ has no poles anywhere in $\C$. The trivial representation is excluded, since its Dedekind zeta function has a pole.
\par\smallskip\textit{Formulation and concept references:} \sref{66}{1}.
\cataloguescope{Prove that every non-trivial irreducible Artin L-function L(s, rho) associated to a Galois representation rho is an entire holomorphic function on the whole complex plane.}
\subsection{Short English statement}
Are all nontrivial irreducible Artin $L$-functions holomorphic throughout the complex plane?
% BEGIN RESEARCH INSERTION
\subsection{Research attempt}

\paragraph{Current frontier.}
Brauer induction provides meromorphic continuation, and one-dimensional Artin functions are controlled by Hecke theory. The missing issue is removal of poles, not continuation itself. Gun--Hazra--Sahu give holomorphy and nonvanishing criteria under specific group and zero-order hypotheses, not unrestricted Artin holomorphy \eref{66}{1}, Introduction and Theorems 1--3. We use only their recalled classical induction identities and analytic facts. The attempt deliberately tries to falsify a positivity lemma that would make induction from smaller groups work.

\paragraph{Attack route.}
Automorphic realization would provide much more than a formal character identity. Brauer induction alone writes a function as a quotient of entire functions and leaves zeros of the denominator uncontrolled. A third route forms a virtual character of orders at a fixed point, then tries to deduce nonnegative coefficients from all proper-subgroup restrictions. We pursue this route and find an explicit obstruction in $A_5$.

\paragraph{Attempt: formal order characters.}
For a finite Galois extension $L/K$ with group $G$, fix $s_0\ne1$ and put
\[
 a_\chi=\operatorname{ord}_{s=s_0}L(s,\chi,L/K),\qquad
 \Theta_{s_0}=\sum_{\chi\in\operatorname{Irr}(G)}a_\chi\chi.
\]
Orders are integers by meromorphy. The desired holomorphy at $s_0$ means $a_\chi\ge0$ for every nontrivial irreducible $\chi$. The trivial coefficient is already nonnegative since the Dedekind zeta function has no pole away from one. For $H\le G$ and a linear character $\lambda$ of $H$, induction gives
\begin{equation}\label{r66:tests}
 \langle\Theta_{s_0},\operatorname{Ind}_H^G\lambda\rangle_G
 =\operatorname{ord}_{s=s_0}L(s,\lambda,L/L^H)\ge0.
\end{equation}
This follows by decomposing the induced representation, taking orders of its finite product, and applying the one-dimensional theorem \eref{66}{1}, Section 2. The tempting intermediate assertion is that these inequalities force a virtual character to be genuine. We disprove that purely character-theoretic assertion.

\paragraph{An explicit virtual character.}
Let $C_5\le A_5$ be generated by a $5$-cycle, and let $A_4$ be a point stabilizer in the natural five-letter action. Set
\begin{equation}\label{r66:theta}
 \Theta=\operatorname{Ind}_{C_5}^{A_5}1-
              \operatorname{Ind}_{A_4}^{A_5}1+1.
\end{equation}
It is a virtual character with integer irreducible coefficients, being a difference of permutation characters. The class values are
\[
\begin{array}{c|rrrrr}
 \text{type}&1&2^2&3&5_a&5_b\\ \hline
 \text{class size}&1&15&20&12&12\\
 \operatorname{Ind}_{C_5}^{A_5}1&12&0&0&2&2\\
 \operatorname{Ind}_{A_4}^{A_5}1&5&1&2&0&0\\
 \Theta&8&0&-1&3&3\\
 \chi_4&4&0&1&-1&-1
\end{array}
\]
Here $\chi_4$ is the actual natural permutation representation modulo its invariant line. Its norm is $(16+20+24)/60=1$, so it is irreducible. The $C_5$ row follows from
\[
 (\operatorname{Ind}_H^G1)(g)
   =\frac{|C_G(g)|}{|H|}\,|H\cap g^G|:
\]
each split $5$-cycle class meets $C_5$ in two elements, and its centralizer has order five. Consequently
\begin{equation}\label{r66:negative}
 \langle\Theta,\chi_4\rangle
   =\frac{8\cdot4+20(-1)\cdot1+24\cdot3(-1)}{60}=-1.
\end{equation}
Thus $\Theta$ is not a genuine character.

\paragraph{Nevertheless every proper restriction is genuine.}
The proper subgroup types are
\[
 1,\ C_2,\ C_3,\ C_5,\ V_4,\ S_3,\ D_5,\ A_4,
\]
with $|D_5|=10$. One can justify this list without relying on numerical sampling. The class sizes show simplicity: no proper union of nonidentity classes, with the identity adjoined, has a nontrivial order dividing $60$. Coset actions then exclude subgroups of indices two, three, and four. An order-$12$ subgroup must fix a letter: otherwise the orbit partition is $2+3$, but the even permutations preserving that partition have order only six. It is therefore a point stabilizer. Orders six and ten give normal Sylow subgroups of orders three and five and the indicated dihedral groups; order four gives $V_4$ since there is no element of order four. The remaining proper orders are prime or one.

Write $\omega$ for a nontrivial $C_3$ character, $\xi$ for a generator of the character group of $C_5$, $\tau$ for the two-dimensional irreducible of $S_3$, $\rho,\rho'$ for the two-dimensional irreducibles of $D_5$, and $\alpha,\bar\alpha,\sigma$ for the nontrivial linear and three-dimensional characters of $A_4$. The restrictions are
\[
\begin{array}{c|l}
 H&\operatorname{Res}^{A_5}_H\Theta\\ \hline
 1&8\cdot1\\
 C_2&4\cdot1+4\cdot\mathrm{sgn}\\
 C_3&2\cdot1+3\omega+3\bar\omega\\
 C_5&4\cdot1+\xi+\xi^2+\xi^3+\xi^4\\
 V_4&2\,\mathrm{reg}_{V_4}\\
 S_3&1+\mathrm{sgn}+3\tau\\
 D_5&2\cdot1+2\cdot\mathrm{sgn}+\rho+\rho'\\
 A_4&\alpha+\bar\alpha+2\sigma
\end{array}
\]
These are exact identities of characters. On $A_4$, the last expression has values $8,0,-1$ on the identity, double transpositions, and $3$-cycles. On $D_5$, $\rho+\rho'$ equals $-1$ at every nonidentity rotation and zero at reflections; adding the linear terms gives $8,3,0$. The cyclic cases follow by summing roots of unity. All coefficients are nonnegative integers.

Finally, $A_5$ is perfect, its only linear character is trivial, and
\[
 \langle\Theta,1\rangle=(8-20+72)/60=1.
\]
Frobenius reciprocity therefore proves
\begin{equation}\label{r66:alltests}
 \langle\Theta,\operatorname{Ind}_H^{A_5}\lambda\rangle\ge0
 \quad(H\le A_5,\ \lambda\text{ linear}).
\end{equation}
More strongly, the inequality holds for every genuine character of every proper subgroup. Hence even positivity of all proper restrictions cannot force positivity of all global irreducible coefficients.

\paragraph{Stress test: this is not an Artin counterexample.}
This refutes only the abstract cone lemma. An actual order character must arise from one extension and one complex point, with additional constraints from Euler products and functional equations. Nothing here realizes~\eqref{r66:theta} as an order character or exhibits an Artin $L$-function with a pole.

Excluding $s_0=1$ removes the trivial zeta pole from the tests; handling that point cannot repair the false implication elsewhere. Negative character values are not pole evidence: the irreducible coefficients, not pointwise values, encode orders.

The accompanying script constructs all $60$ even permutations and all $59$ subgroups. It computes induced permutation values and the negative inner product exactly, and checks nonnegativity of all induced-linear tests by Fourier analysis on the finite abelianizations using rational matrices: all 159 linear-character tests pass, with minimum multiplicity zero. These checks concern one finite group, not zeros or analytic continuation. The displayed decompositions supply an independent hand-verifiable check of every proper restriction.

\paragraph{Outcome.}
\textbf{Outcome: Exploratory approach with unresolved gap.}

The attempt supplies a concrete counterexample to a plausible intermediate positivity statement, even after strengthening its hypotheses to genuine restrictions on every proper subgroup. This blocks that induction route, not Artin's conjecture. Continuing would require an additional analytic constraint excluding virtual characters of this kind. No such general constraint is proved; no Artin counterexample or priority for the finite-character construction is claimed.

% END RESEARCH INSERTION
\subsection{Sources}
\begin{itemize}\raggedright
\item[\textnormal{[S1]}]\hypertarget{source:66:1}{} E. Artin, Zur Theorie der L-Reihen mit allgemeinen Gruppencharakteren, Abh. Math. Sem. Univ. Hamburg 8:292-306, 1930.
\url{https://doi.org/10.1073/pnas.16.8.532}.
{\small\textit{Catalogue formulation source.}}
\item[\textnormal{[S2]}]\hypertarget{source:66:2}{} On holomorphy and non-vanishing of Artin L-functions.
\url{https://link.springer.com/article/10.1007/s00605-025-02060-7}.
{\small\textit{Catalogue research/status reference; not a proof endorsement.}}
\item[\textnormal{[S3]}]\hypertarget{source:66:3}{} Artin's Conjectures — Ram Murty, Fields Institute, 12 March 2025.
\url{https://www1.fields.utoronto.ca/talks/Artins-Conjectures}.
{\small\textit{Catalogue research/status reference; not a proof endorsement.}}

% BEGIN ADDED RESEARCH REFERENCES
\item[\textnormal{[E1]}]\hypertarget{extra:66:1}{} Sanoli Gun, Suhita Hazra, and Dhananjaya Sahu, \emph{On holomorphy and non-vanishing of Artin $L$-functions}, Monatshefte für Mathematik 206 (2025), 853--883, DOI 10.1007/s00605-025-02060-7; Introduction and Section 2.
\url{https://link.springer.com/article/10.1007/s00605-025-02060-7}.
{\small\textit{Classical Artin induction/Hecke facts and the scope of the paper's conditional and special-group holomorphy results. Consulted 27 September 2026.}}
% END ADDED RESEARCH REFERENCES
\end{itemize}

\end{document}
